\documentclass[11pt, a4paper]{article}

% --- PACCHETTI ---
\usepackage[utf8]{inputenc}
\usepackage[italian]{babel}
\usepackage[margin=2.5cm]{geometry} % Imposta i margini della pagina
\usepackage{graphicx}               % Per inserire il logo
\usepackage{tabularx}               % Per tabelle larghe quanto la pagina
\usepackage[table]{xcolor}          % Per colorare le celle delle tabelle
\usepackage[hidelinks]{hyperref}    % Per stile link cliccabili
\usepackage{lmodern}                % Rende i font scalabili
\usepackage{anyfontsize}            % Permette dimensioni personalizzate

\begin{document}

% ==============================================================================
% FRONTESPIZIO (PAGINA 1)
% ==============================================================================
\begin{titlepage}
    \centering
    \vspace*{0.5cm}

    % LOGHI
    \includegraphics[width=0.45\textwidth]{assets/logo_unipd.jpg}
    \hfill
    \includegraphics[width=0.45\textwidth]{assets/logo.png}\par
    \vspace{0.5cm}

    % NOME DEL GRUPPO
    {\fontsize{45}{50}\sffamily\selectfont \textbf{On-By-One}\par} 
    \vspace{1cm}

    % TITOLO DEL VERBALE
    {\huge \bfseries Verbale interno 2026/11/10 \par}
    
    \vspace{1.5cm}

    % MEMBRI DEL GRUPPO
    {\large
    \textbf{On-By-One}:\par
    \vspace{0.3cm}
    \begin{tabular}{l l}
        Sara Ballardin & 2101048 \\
        Shuhui Chen & 2138001 \\
        Simone Giuliano & 2145591 \\
        Alvise Greggio & 2141004 \\
        Daniele Nuzzaci & 2139039 \\
        Federico Temporin & 2146764 \\
        Tommaso Zorzetto & 2148629 \\
    \end{tabular}
    \par}

    \vfill

    % TABELLA INFORMAZIONI DOCUMENTO E DESTINATARI
    \begin{table}[h]
        \centering
        \arrayrulecolor{gray!50}
        \renewcommand{\arraystretch}{1.5}
        \begin{tabularx}{\textwidth}{|c|c|c|X|}
            \hline
            \rowcolor{gray!10} \multicolumn{4}{|c|}{\textbf{Informazioni documento}} \\
            \hline
            \rowcolor{gray!10} \textbf{Versione} & \textbf{Data} & \textbf{Stato} & \multicolumn{1}{c|}{\textbf{Destinatari}} \\
            \hline
            1.0.0 & 2026/11/10 & Approvato & 
            On-By-One \newline 
            Prof. Tullio Vardanega \newline 
            Prof. Riccardo Cardin \\
            \hline
        \end{tabularx}
        \arrayrulecolor{black}
    \end{table}
    
    \vspace{0.5cm}
    
    % CONTATTI
    {\large Contatti: \href{mailto:onbyone.swe@gmail.com}{onbyone.swe@gmail.com} }
    
\end{titlepage}

\newpage

% ==============================================================================
% REGISTRO DELLE MODIFICHE (PAGINA 2)
% ==============================================================================
\section*{Registro delle modifiche}

\vspace{1cm}
\begin{center}
    \arrayrulecolor{gray!50}
    \renewcommand{\arraystretch}{1.5}
    \begin{tabularx}{\textwidth}{|c|c|c|X|c|}
        \hline
        \rowcolor{gray!10} 
        \multicolumn{1}{|c|}{\textbf{Versione}} & 
        \multicolumn{1}{c|}{\textbf{Data}} & 
        \multicolumn{1}{c|}{\textbf{Autore}} & 
        \multicolumn{1}{c|}{\textbf{Descrizione}} & 
        \multicolumn{1}{c|}{\textbf{Verificatore}} \\
        \hline
        1.0.0 & 2026/11/10 & Tommaso Zorzetto & Stesura del verbale & Daniele Nuzzaci \\
        \hline
        0.1.0 & 2026/11/10 & Tommaso Zorzetto & Stesura bozza iniziale & - \\
        \hline
    \end{tabularx}
    \arrayrulecolor{black}
\end{center}
\vfill
\newpage

% ==============================================================================
% INDICE (PAGINA 3)
% ==============================================================================
\tableofcontents
\newpage

% ==============================================================================
% INFORMAZIONI GENERALI
% ==============================================================================
\section{Informazioni generali}

\begin{table}[h]
    \centering
    \arrayrulecolor{gray!50}
    \renewcommand{\arraystretch}{1.5}
    \begin{tabularx}{\textwidth}{|c|>{\centering\arraybackslash}X|}
        \hline
        \rowcolor{gray!10} \multicolumn{2}{|c|}{\textbf{Dettagli della riunione}} \\
        \hline
        \textbf{Data} & 2026/11/10 \\ \hline
        \textbf{Ora inizio} & 15:00 \\ \hline
        \textbf{Ora fine} & 17:00 \\ \hline
        \textbf{Modalità} & Discord \\ \hline
    \end{tabularx}
    \arrayrulecolor{black}
\end{table}

\begin{table}[h]
    \centering
    \arrayrulecolor{gray!50}
    \renewcommand{\arraystretch}{1.5}
    \begin{tabularx}{\textwidth}{|>{\centering\arraybackslash}X|>{\centering\arraybackslash}X|>{\centering\arraybackslash}X|}
        \hline
        \rowcolor{gray!10} \multicolumn{3}{|c|}{\textbf{Partecipanti}} \\
        \hline
        \rowcolor{gray!10} \textbf{Nome e cognome} & \textbf{Ruolo} & \textbf{Presenza} \\
        \hline
        Sara Ballardin & Analista & Presente \\ \hline
        Shuhui Chen & Analista & Presente \\ \hline
        Simone Giuliano & Analista & Presente \\ \hline
        Alvise Greggio & Responsabile & Presente \\ \hline
        Daniele Nuzzaci & Verificatore & Presente \\ \hline
        Federico Temporin & Amministratore & Presente \\ \hline
        Tommaso Zorzetto & Verificatore & Presente \\ \hline
    \end{tabularx}
    \arrayrulecolor{black}
\end{table}

% ==============================================================================
% CORPO DEL VERBALE
% ==============================================================================

\section{Ordine del giorno}
\begin{itemize}
    \item Designazione del redattore e preparazione delle domande tecniche per il primo incontro in presenza con l'azienda proponente Vimar.
    \item Chiarimento della ripartizione del carico di lavoro e dei criteri di rotazione dei ruoli all'interno del team.
    \item Strategia di redazione del documento \textit{Analisi dei Requisiti} (gerarchia tra requisiti e casi d'uso, divisione delle sezioni del capitolato).
    \item Definizione del ruolo e delle mansioni del verificatore per la valutazione delle scelte tecnologiche e dei documenti.
    \item Impostazione operativa del \textit{Piano di Progetto} e del \textit{Piano di Qualifica} (rendicontazione oraria, consuntivo di periodo e diagrammi di avanzamento).
    \item Organizzazione delle \textit{milestone}, delle etichette (label) e della struttura della repository per la revisione \textbf{RTB}.
\end{itemize}

\newpage

\section{Discussioni}

\subsection{Preparazione all'incontro con l'azienda Vimar e designazione redattore}
In apertura di riunione il gruppo ha designato Simone Giuliano come redattore principale incaricato di presenziare e verbalizzare l'incontro presso l'azienda proponente Vimar, concordando di richiedere preventivamente l'autorizzazione alla registrazione audio della parte dedicata alle domande tecniche. Il team ha elaborato i quesiti da porre ai referenti aziendali: chiarire se le tecnologie suggerite nel capitolato rappresentino vincoli stringenti o alternative aperte a valutazione, e raccogliere indicazioni pratiche sul livello di granularità e sul punto di partenza dell'analisi (se procedere \textit{top-down} dai macro-requisiti ai casi d'uso o viceversa).

\subsection{Gestione dei ruoli, rotazione e compiti di verifica}
Il team si è confrontato sulla gestione dei ruoli in funzione delle direttive didattiche. Riguardo alle mansioni di verifica, è stato stabilito che il verificatore non deve duplicare l'indagine svolta dall'analista, bensì accertare la fondatezza logica e metodologica delle motivazioni addotte, la coerenza formale e l'assenza di contraddizioni macroscopiche.

\subsection{Pianificazione e redazione dell'Analisi dei Requisiti}
È stata discussa la struttura redazionale dell'\textit{Analisi dei Requisiti}. Il gruppo ha convenuto di strutturare il testo partendo dai requisiti (funzionali, qualitativi e di vincolo) estratti direttamente dal capitolato d'appalto per poi derivare in modo coerente i casi d'uso. Gli analisti (Sara Ballardin, Shuhui Chen, Simone Giuliano) hanno pianificato la ripartizione del testo del capitolato (paragrafi 2 e 3) e concordato una scadenza intermedia entro il fine settimana per la stesura della prima bozza strutturata.

\subsection{Impostazione del Piano di Progetto e del Piano di Qualifica}
Si è esaminata l'organizzazione dei documenti di pianificazione e controllo qualità. Per il \textit{Piano di Progetto}, Alvise Greggio recepirà la struttura di base predisposta da Sara Ballardin per completare le sezioni relative all'analisi dei rischi, ai consuntivi di periodo dello sprint e alla pianificazione temporale (valutando l'integrazione di diagrammi Gantt e PERT). Per il \textit{Piano di Qualifica}, si è rimandata la stesura analitica alla trattazione approfondita delle metriche e del cruscotto di valutazione, concentrando per ora l'attenzione sui macro-obiettivi di qualità di processo e prodotto.

\subsection{Organizzazione della repository, etichette e milestone per RTB}
Federico Temporin ha coordinato la configurazione dell'ambiente di lavoro su GitHub. È stata creata la cartella \texttt{docs/tex/rtb/} per accogliere i documenti specifici per la prima revisione di avanzamento. Sono state inoltre definite tre milestone principali associate ai documenti cardine (\textit{Analisi dei Requisiti}, \textit{Piano di Progetto}, \textit{Piano di Qualifica}), ed è stata introdotta l'etichetta \texttt{sprint-1} per tracciare puntualmente le issue appartenenti all'iterazione corrente.

\newpage

\section{Decisioni}
\begin{itemize}
    \item \texttt{VI-2026-11-10.1} Assegnato a Simone Giuliano il ruolo di redattore per l'incontro in presenza con l'azienda Vimar, con l'indicazione di richiedere il consenso alla registrazione per la sessione di domande tecniche.
    \item \texttt{VI-2026-11-10.2} Approvata la formulazione dei quesiti tecnici per Vimar relativi al grado di vincolo delle tecnologie suggerite e alle indicazioni operative sull'architettura.
    \item \texttt{VI-2026-11-10.3} Definito il criterio di lavoro per l'\textit{Analisi dei Requisiti}: estrazione iniziale dei requisiti espliciti dal capitolato e successiva definizione dei casi d'uso correlati.
    \item \texttt{VI-2026-11-10.4} Deliberata la creazione della directory \texttt{docs/rtb/} e l'introduzione delle milestone RTB per i tre documenti cardine (\textit{Analisi dei Requisiti}, \textit{Piano di Progetto} e \textit{Piano di Qualifica}).
    \item \texttt{VI-2026-11-10.5} Adottata l'etichetta \texttt{sprint-1} per l'aggregazione e il tracciamento granulare di tutte le issue appartenenti al primo sprint.
\end{itemize}

\section{To-do}
\begin{table}[h]
    \centering
    \arrayrulecolor{gray!50}
    \renewcommand{\arraystretch}{1.5}
    \begin{tabularx}{\textwidth}{|c|X|p{3.15cm}|c|}
        \hline
        \rowcolor{gray!10} 
        \multicolumn{1}{|c|}{\textbf{ID issue}} & 
        \multicolumn{1}{c|}{\textbf{Task}} & 
        \multicolumn{1}{c|}{\textbf{Assegnatario}} & 
        \multicolumn{1}{c|}{\textbf{Scadenza}} \\
        \hline
        \href{https://github.com/On-By-One-Team/On-By-One.github.io/issues/68}{\#68} & Stesura bozza del documento Piano di Progetto & Sara Ballardin & 2026/11/15 \\
        \hline
        \href{https://github.com/On-By-One-Team/On-By-One.github.io/issues/69}{\#69} & Nuova organizzazione per RTB & Federico Temporin & 2026/11/13 \\
        \hline
        \href{https://github.com/On-By-One-Team/On-By-One.github.io/issues/70}{\#70} & Stesura del verbale interno del 2026/11/10 & Tommaso Zorzetto & 2026/11/12 \\
        \hline
        - & Stesura della bozza iniziale di Analisi dei Requisiti & Sara Ballardin, Shuhui Chen, Simone Giuliano & 2026/11/15 \\
        \hline
        - & Avanzamento del Piano di Progetto (rischi, rendicontazione e diagrammi) & Alvise Greggio & 2026/11/16 \\
        \hline
        - & Aggiornamento dei Diari di Bordo & Alvise Greggio & 2026/11/11 \\
        \hline
    \end{tabularx}
    \arrayrulecolor{black}
\end{table}

\end{document}
